\subsubsection{开源代码逻辑、约束与选型边界}

代码流水线为：
\begin{quote}
\small\ttfamily
schema -> audit -> split -> preprocess -> baseline\\
-> fit -> validate -> explain -> export
\end{quote}
Python 侧使用 pandas、scikit-learn、
statsmodels 及主流梯度提升输入输出关系；MATLAB 侧建立读表、缺失处理、回归分类和时序
对照。所有随机切分记录种子，时间序列禁止随机打乱，分组样本按主体隔离。
评价模型在正向化、无量纲化和权重确定后再计算排序，并做权重扰动和秩相关检验。

K-Means 假定近似球形且需预设簇数，DBSCAN 能识别噪声和任意形状但依赖密度尺度；
层次聚类便于解释树状结构但大样本代价高；谱聚类适合非凸簇但需构图。Lasso 可
稀疏选择，Ridge 对共线性更稳定；SVR 适合中小样本但调参和扩展性受限；随机森林
稳健易用但外推弱；梯度提升精度高但要防泄漏和过拟合。ARIMA/SARIMA 要检验
平稳和季节性，VAR 要控制维数和滞后，GM(1,1) 只适合小样本单调趋势，深度时序
模型不能替代滚动窗口基线。

